\documentclass[aspectratio=169,10pt]{beamer}
\usetheme{Madrid}
\usecolortheme{default}

\usepackage{hyperref}
\usepackage{tikz}
\usepackage{amsmath,amssymb}

% --- Custom Colors ---
\definecolor{unibblue}{RGB}{0,51,102}
\definecolor{unibgold}{RGB}{212,175,55}
\definecolor{unibgreen}{RGB}{0,128,0}

\setbeamercolor{palette primary}{bg=unibblue,fg=white}
\setbeamercolor{palette secondary}{bg=unibblue!85,fg=white}
\setbeamercolor{palette tertiary}{bg=unibblue!70,fg=white}
\setbeamercolor{structure}{fg=unibblue}
\setbeamercolor{block title}{bg=unibblue!85,fg=white}
\setbeamercolor{block body}{bg=unibblue!10,fg=black}
\setbeamercolor{alertblock title}{bg=unibgold!90,fg=black}
\setbeamercolor{alertblock body}{bg=unibgold!15,fg=black}

\title{Persamaan Diferensial}
\subtitle{Minggu 14: Fungsi Khusus — Persamaan Bessel \& Aplikasi Gelombang Silinder}
\author{Novalio Daratha \& Muhammad Arfan}
\institute{Program Studi Teknik Elektro\\Universitas Bengkulu}
\date{2026}

\begin{document}

\begin{frame}
  \titlepage
\end{frame}

\begin{frame}{Capaian Pembelajaran (CPMK 4)}
  \begin{itemize}
    \item \textbf{CPMK 4.1:} Memahami transformasi persamaan diferensial dari koordinat Kartesius $(x,y)$ ke koordinat Silinder $(r,\theta,z)$.
    \item \textbf{CPMK 4.2:} Menyusun bentuk standar Persamaan Diferensial Bessel dan mengenali solusi fungsi Bessel Jenis Pertama $J_\nu(x)$ dan Jenis Kedua $Y_\nu(x)$.
    \item \textbf{CPMK 4.3:} Memahami perilaku nilai nol (*zeros*) dari fungsi Bessel untuk penentuan frekuensi resonansi rongga silinder.
    \item \textbf{CPMK 4.4:} Menganalisis fenomena Efek Kulit (*Skin Effect*) dan distribusi medan pada kabel koaksial serta pandu gelombang (*waveguide*).
  \end{itemize}
\end{frame}

\begin{frame}{Daftar Isi}
  \tableofcontents
\end{frame}

\section{Asal Mula Persamaan Bessel dalam Koordinat Silinder}
\begin{frame}{Laplacian dalam Koordinat Silinder}
  Pada sistem berpenampang melingkar (kabel tembaga silinder, serat optik, *cavity resonator*):
  $$ \nabla^2 u = \frac{1}{r}\frac{\partial}{\partial r}\left(r\frac{\partial u}{\partial r}\right) + \frac{1}{r^2}\frac{\partial^2 u}{\partial \theta^2} + \frac{\partial^2 u}{\partial z^2} = 0 $$
  \pause
  Dengan pemisahan variabel radial $R(r)$, diperoleh:
  \begin{alertblock}{Bentuk Standar Persamaan Diferensial Bessel Orde $\nu$}
    $$ {\color{unibblue} x^2 \frac{d^2y}{dx^2} + x \frac{dy}{dx} + (x^2 - \nu^2)y = 0} $$
    dengan parameter $\nu \ge 0$ (orde Bessel).
  \end{alertblock}
\end{frame}

\section{Solusi: Fungsi Bessel $J_\nu(x)$ \& $Y_\nu(x)$}
\begin{frame}{Fungsi Bessel Jenis Pertama \& Kedua}
  \begin{columns}
    \begin{column}{0.5\textwidth}
      \begin{block}{Fungsi Bessel Jenis 1: $J_\nu(x)$}
        \begin{itemize}
          \item Berhingga di titik pusat $x = 0$ ($J_0(0) = 1, J_n(0) = 0$ untuk $n \ge 1$).
          \item Memodelkan medan elektromagnetik di dalam silinder pejal.
        \end{itemize}
      \end{block}
    \end{column}
    \begin{column}{0.5\textwidth}
      \pause
      \begin{block}{Fungsi Bessel Jenis 2: $Y_\nu(x)$ (Neumann)}
        \begin{itemize}
          \item Menuju $-\infty$ saat $x \to 0$ (singular di pusat).
          \item Dipakai untuk wilayah cincin/anulus (misal antara konduktor dalam dan luar kabel koaksial).
        \end{itemize}
      \end{block}
    \end{column}
  \end{columns}
  \pause
  \begin{alertblock}{Solusi Umum Persamaan Bessel}
    $$ y(x) = C_1 J_\nu(x) + C_2 Y_\nu(x) $$
  \end{alertblock}
\end{frame}

\section{Aplikasi Elektro: Efek Kulit \& Pandu Gelombang}
\begin{frame}{Aplikasi 1: Fenomena Efek Kulit (*Skin Effect*)}
  \begin{block}{Distribusi Rapat Arus Frekuensi Tinggi}
    Pada frekuensi tinggi ($f \ge 50\text{ Hz}$ AC), arus listrik terdesak ke permukaan luar konduktor:
    $$ J_z(r) = J_0 \frac{J_0(k r)}{J_0(k a)}, \quad k = \frac{1 - j}{\delta} $$
    dengan kedalaman kulit (*skin depth*) $\delta = \sqrt{\frac{1}{\pi f \mu \sigma}}$.
  \end{block}
  \pause
  \textit{Dampak Rekayasa: Menjelaskan mengapa saluran transmisi tegangan ekstra tinggi (SUTET) menggunakan kabel konduktor berongga atau berkas (*bundled conductor*).}
\end{frame}

\section{Kuis \& Rangkuman}
\begin{frame}{Kuis Interaktif 14}
  \textbf{Soal:} Mengapa pada analisis kabel silinder tembaga pejal, konstanta $C_2$ pada suku $Y_\nu(x)$ harus selalu bernilai nol ($C_2 = 0$)?
  \pause
  \begin{block}{Jawaban:}
    Karena di titik pusat kabel ($r = 0$), nilai fungsi Bessel jenis kedua $Y_\nu(0) \to -\infty$.
    Secara fisik, rapat arus atau medan listrik di sumbu kabel tidak mungkin bernilai minus tak hingga (harus bernilai terhingga dan riil). Maka suku $Y_\nu$ harus dihilangkan dengan memilih $C_2 = 0$.
  \end{block}
\end{frame}

\begin{frame}{Rangkuman Materi Minggu 14}
  \begin{enumerate}
    \item Koordinat silinder secara alami memunculkan PDB Bessel untuk bagian radial.
    \item $J_\nu(x)$ berosilasi mirip fungsi kosinus/sinus tetapi amplitudonya meluruh seiring $\frac{1}{\sqrt{x}}$.
    \item Fungsi Bessel adalah pilar komputasi elektromagnetika pada pandu gelombang silinder dan serat optik.
  \end{enumerate}
\end{frame}

\end{document}
